% ===========================================================================
\section{The norm bound: a sum-of-squares zerocheck}\label{sec:norm}
% ===========================================================================

The acceptance bound \eqref{eq:norm} is a \emph{direct} sum of squares
over $\Z$ of the committed coefficients --- not a per-row accumulator
recurrence. We introduce a non-negative slack and arithmetize, per
signature $j$,
\begin{equation}\label{eq:normeq}
  \sum_{i=0}^{n-1} s_1[i]^2 + \sum_{i=0}^{n-1} s_2[i]^2 + \col{slack}(j)
  \;=\; \nbound, \qquad \col{slack}(j) \ge 0 .
\end{equation}

\paragraph{Coefficient slices.}
The norm needs the individual integer coefficients $s_1[i], s_2[i]$.
These are already present, uncombined, as the coefficients of the limb
cells $\col{s_1}_m, \col{s_2}_m$. We read them with the
\emph{coefficient-slice} machinery --- the integer-coefficient analogue
of the booleanity check's bit-slice extraction. For a selected set of
limb columns, \texttt{compute\_coeff\_slices\_flat} produces, for each
coefficient position, an MLE over the signature rows whose value at row
$j$ is that coefficient of that limb in signature $j$. The norm group
consumes the $2N$ slices of the $\col{s_1},\col{s_2}$ limb cells (here
$N = n = LW$), plus the per-row $\col{slack}$.

\paragraph{The zerocheck.}
Following \texttt{piop/src/lookup/booleanity.rs}, the bound is enforced as
an $\eqop$-weighted zerocheck over the $2^{\nu}$ signature rows. With
$r$ the sumcheck challenge point and $\{\mathrm{slice}_k\}_{k<2N}$ the
coefficient slices, the group proves
\begin{equation}\label{eq:norm-zerocheck}
  \sum_{j \in \{0,1\}^{\nu}} \eqop(r, j)\,
  \Bigl( \underbrace{\textstyle\sum_{k=0}^{2N-1} \mathrm{slice}_k(j)^2
         + \col{slack}(j)}_{=:\ \mathrm{comb}(j)}
         \;-\; \nbound \Bigr) \;=\; 0 .
\end{equation}
The per-row combiner $\mathrm{comb}(j) = \sum_k \mathrm{slice}_k(j)^2 +
\col{slack}(j)$ is \texttt{falcon\_norm\_comb\_fn}; the group multiplies
by $\eqop(r,j)$ and subtracts the public constant $\nbound$. Because the
$\eqop$ weights select each row independently, \eqref{eq:norm-zerocheck}
holds iff \eqref{eq:normeq} holds for \emph{every} signature --- exactly as
booleanity's $\sum_j \eqop(r,j)(v_j^2 - v_j) = 0$ enforces $v_j\in\{0,1\}$
per row.

\paragraph{Degree and group.}
The combiner is degree~$2$ in the slices (the squares); with the
$\eqop$ factor the group is degree~$3$. Unlike booleanity, the
coefficients are not bits, so there is no round-one fast path: the group
uses the plain multi-degree sumcheck constructor
(\texttt{MultiDegreeSumcheckGroup::new}) with claimed sum~$0$. It is run
\emph{alongside} the combined-polynomial resolver group in the same
multi-degree sumcheck (group order $[\,\text{cpr}, (\text{booleanity?}),
\text{norm}\,]$; Falcon has no binary columns, so booleanity is absent and
the norm group sits at index~$1$).

\paragraph{Binding to the committed columns.}
The verifier binds the slice evaluations to the committed
$\col{s_1},\col{s_2}$ limb cells by the projection-element consistency
check
$\mle{\col{P}_m}(r^\ast) = \sum_{p<W} a^{\,p}\,\mathrm{slice}_{mW+p}(r^\ast)$,
where $a$ is the evaluation-projection element ($\evalproj{a}$) used in
the inner field PIOP, and binds $\col{slack}$ via its own committed-column
opening. No separate equality argument is needed.

\paragraph{Range checks (out of scope here).}
Two non-negativity / range facts complete soundness and are a separate
bit-decomposition $+$ booleanity pass, not yet wired in the reference
instance: \emph{(i)}~$\col{slack}(j)\ge 0$ (so \eqref{eq:normeq} is a
genuine upper bound, not an equality that a large $\col{slack}$ could
satisfy from above); and \emph{(ii)}~the centering
$s_1[i],s_2[i]\in(-\half,\half\,]$, which pins the representatives in
\eqref{eq:s1}. Both are standard range arguments over the same limb cells.
